\chapter{Introducción}

\section{Contexto General y Planteamiento del Problema}

La electroencefalografía (EEG) ha sido durante décadas una de las herramientas más importantes en el ámbito de la neurociencia clínica y experimental. Gracias a su alta resolución temporal, su carácter no invasivo y su costo relativamente bajo en comparación con otras técnicas de neuroimagen —como la Imagen por Resonancia Magnética Funcional (fMRI) o la Tomografía por Emisión de Positrones (PET)—, el EEG se ha convertido en el estándar de oro para el estudio de la dinámica cerebral en tiempo real. Sus aplicaciones abarcan desde el diagnóstico clínico de patologías neurológicas (como la epilepsia y los trastornos del sueño) hasta la investigación en neurociencia cognitiva y el desarrollo avanzado de interfaces cerebro-computador (BCI, por sus siglas en inglés).

Sin embargo, a pesar de sus vastas ventajas, la adquisición de señales EEG es un proceso sumamente delicado y susceptible a múltiples fuentes de contaminación. En la práctica clínica y de laboratorio, es habitual enfrentar la degradación o pérdida total de información en uno o más electrodos (canales). Estas fallas pueden originarse por diversas razones técnicas y ambientales: aumento en la impedancia de contacto entre el electrodo y el cuero cabelludo debido a la transpiración o secado del gel conductor, fallas de hardware en el amplificador (saturación de voltaje), movimientos bruscos del paciente que desconectan físicamente los sensores, o simplemente el descarte deliberado de canales durante la fase de preprocesamiento debido a una relación señal-ruido inaceptablemente baja.

La pérdida de canales no es un inconveniente menor; representa una amenaza crítica para la integridad del análisis posterior. La ausencia de información espacial interrumpe la topografía de la señal, distorsionando la localización de fuentes neuronales (source localization) y alterando biomarcadores críticos como los Potenciales Relacionados a Eventos (ERPs) y la Densidad Espectral de Potencia (PSD). Además, en el contexto de los modernos sistemas de \textit{Machine Learning} aplicados a BCIs, la pérdida impredecible de canales interrumpe la dimensionalidad de entrada de los clasificadores, obligando frecuentemente a descartar ensayos completos (trials) de datos valiosos. En consecuencia, la interpolación precisa de estos canales perdidos se vuelve un paso ineludible en el \textit{pipeline} de preprocesamiento del EEG.

\section{Limitaciones de los Enfoques Tradicionales}

Históricamente, el problema de la interpolación de canales EEG ha sido abordado mediante métodos espaciales o estadísticos. El enfoque más extendido y utilizado por defecto en \textit{toolboxes} de análisis como EEGLAB o MNE-Python es la interpolación por \textit{Spherical Splines} \cite{perrin1989}. Este método, junto con otros enfoques basados en la distancia euclidiana (como la interpolación por vecinos más cercanos o funciones de base radial), asume que la cabeza humana es una esfera perfecta y estima el voltaje del canal faltante proyectando geométricamente los voltajes de los electrodos circundantes.

Aunque estos métodos son rápidos y matemáticamente elegantes, adolecen de una limitación fundamental: operan estrictamente en el dominio espacial estático para cada instante de tiempo individual. Al hacerlo, ignoran dos realidades fundamentales de la neurofisiología. En primer lugar, la conectividad funcional real del cerebro no es isotrópica: la actividad eléctrica sigue rutas anatómicas complejas que no siempre se correlacionan con la distancia física en el cuero cabelludo, de modo que un modelo esférico resulta una simplificación excesiva. En segundo lugar, las señales EEG son generadas por osciladores neuronales sujetos a leyes físicas de continuidad y momento. Al reconstruir un canal muestra por muestra de forma aislada, los métodos puramente espaciales son incapaces de preservar la dinámica de alta frecuencia, generando trazos ruidosos o distorsiones de fase.

Por otro lado, métodos estadísticos como el Análisis de Componentes Independientes (ICA) o el Análisis de Componentes Principales (PCA) intentan proyectar las señales mixtas en fuentes subyacentes latentes, reconstruyendo el canal perdido a partir de dicha proyección. Sin embargo, en escenarios de pérdida severa (por ejemplo, múltiples canales espacialmente contiguos fallando al mismo tiempo), estos métodos tienden a inyectar ruido global en el canal reconstruido o a fallar en la captura de picos transitorios locales muy rápidos.

\section{El Advenimiento del Procesamiento de Señales en Grafos (GSP)}

Para superar estas limitaciones, en los últimos años ha emergido con fuerza el marco del \textbf{Procesamiento de Señales en Grafos (GSP, Graph Signal Processing)} \cite{shuman2013,ortega2018}. El GSP ofrece una perspectiva natural y poderosa para lidiar con datos que residen en dominios irregulares no euclidianos. En lugar de modelar la cabeza como un espacio continuo tridimensional, GSP modela la red de electrodos como un grafo discreto $\mathcal{G} = (\mathcal{V}, \mathcal{E}, \mathbf{W})$, donde cada electrodo es un nodo (vértice) y las relaciones de similitud o distancia fisiológica entre ellos se codifican en las aristas (pesos).

Bajo este paradigma, el voltaje medido en un instante de tiempo se considera una ``señal de grafo''. GSP permite traducir conceptos fundamentales del procesamiento de señales clásico (como la Transformada de Fourier, el filtrado paso-bajo y el muestreo) al dominio del grafo a través del operador del Laplaciano del grafo ($\mathbf{L}$). La principal intuición detrás de la interpolación GSP es la \textit{suavidad} (smoothness): se asume que, en la naturaleza, las señales conectadas por aristas de alto peso tienden a tener valores similares. Minimizar una función de costo que penalice las variaciones abruptas a través de la topología del grafo permite reconstruir los datos faltantes con mucha mayor precisión que la simple distancia euclidiana \cite{narang2013}.

Aún más avanzado es el concepto de regularización espacio-temporal. Recientes extensiones de GSP han introducido enfoques como la \textit{Regularización Temporal de Suavidad de Sobolev} (TRSS, \textit{Temporal Regularized Sobolev Smoothness}) \cite{giraldo2022}, que optimiza conjuntamente la topología espacial (el grafo) y la derivada temporal de la señal. Al penalizar saltos bruscos en el tiempo de forma simultánea a los saltos en el espacio, TRSS impone restricciones fisiológicas reales, lo que teóricamente debería traducirse en una morfología de onda mucho más fidedigna.

\section{La Brecha en la Literatura y la Motivación del Trabajo}

A pesar de la inmensa promesa teórica de los algoritmos GSP y TRSS, su validación empírica en la literatura de señales biomédicas sufre de una preocupante fragmentación metodológica. En concreto, la literatura presenta dos problemas graves de benchmarking. Por un lado, es habitual que los estudios que proponen nuevos algoritmos GSP se comparen contra \textit{Spherical Splines} u otros métodos clásicos utilizando las configuraciones por defecto de los paquetes de software, sin realizar una búsqueda rigurosa de hiperparámetros para los \textit{baselines}; esto infla artificialmente las ventajas del nuevo método propuesto. Por otro lado, la gran mayoría de los trabajos restringen su evaluación a métricas de error de amplitud global (como el RMSE o MAE) y aplican escenarios de desconexión irrealistas, como la remoción de canales aleatorios aislados en bases de datos pequeñas. Muy rara vez se investiga el impacto morfológico de la reconstrucción mediante métricas temporales como el \textit{Dynamic Time Warping} (DTW) o en el dominio de las frecuencias (LSD), los cuales constituyen los verdaderos predictores de la utilidad clínica de la señal interpolada.

La motivación concreta de este trabajo nace directamente de estas brechas. El objetivo no es diseñar desde cero un nuevo marco teórico matemático, sino someter a la frontera del conocimiento actual (métodos geométricos, GSP y regularización temporal) a la evaluación empírica más rigurosa, exhaustiva y justa posible. Al utilizar frameworks modernos de optimización bayesiana (como \textbf{Optuna} \cite{akiba2019optuna}) para llevar todos los algoritmos a su máximo límite empírico bajo el protocolo evaluado en múltiples datasets, esta tesis busca revelar empíricamente el verdadero ``cruce de rendimiento'' entre el paradigma clásico y el paradigma de grafos.

\section{Objetivos}

\subsection{Objetivo General}
Comparar el rendimiento espacial, temporal y espectral de los métodos de reconstrucción de canales faltantes de EEG basados en Procesamiento de Señales en Grafos (GSP) y regularización temporal frente a baselines geométricos clásicos, utilizando un marco experimental riguroso, reproducible y exhaustivamente optimizado.

\subsection{Objetivos Específicos}
\begin{itemize}
    \item \textbf{Consolidación de Métodos:} Implementar y unificar una amplia colección de métodos de interpolación (geométricos, estadísticos y basados en grafos) en un \textit{pipeline} modular en Python.
    \item \textbf{Diseño de Escenarios Clínicos:} Definir protocolos realistas de simulación de canales faltantes, abarcando desde fallos esporádicos (\textit{random}) hasta pérdidas masivas contiguas de amplificadores locales (\textit{nearby}), iterando desde un $10\%$ hasta un $40\%$ de degradación de la malla sensora.
    \item \textbf{Optimización Imparcial (Fair Benchmarking):} Integrar el \textit{framework} de optimización de hiperparámetros Optuna para explorar iterativamente los espacios paramétricos de todos los métodos (incluyendo las constantes de regularización $\alpha$ y $\beta$, los radios de influencia $\sigma$, y la esparsidad $k$), garantizando que ningún algoritmo sea evaluado por debajo de su límite de diseño.
    \item \textbf{Evaluación Multi-Métrica Transversal:} Cuantificar el desempeño de reconstrucción cruzando tres bases de datos heterogéneas (PhysioNet de alta densidad, BCI IV 2a de baja densidad y MNE Sample para evocados), utilizando una batería robusta de métricas que incluyan no solo error absoluto (MAE, RMSE, SNR), sino también fidelidad morfológica y espectral (DTW, Log Spectral Distance y Coherencia Magnitud-Cuadrada).
    \item \textbf{Análisis de Relevancia Fisiológica:} Demostrar visual y cualitativamente el impacto de los distintos enfoques de interpolación en la preservación de Potenciales Relacionados a Eventos (ERPs) específicos (e.g., N100, P300) y de la dinámica de bandas de frecuencia naturales del cerebro.
    \item \textbf{Trazabilidad Científica:} Generar un repositorio abierto de artefactos computacionales y bitácoras de resultados que permita la reproducibilidad total de los hallazgos para respaldar la redacción del documento de tesis y artículos científicos asociados.
\end{itemize}

\section{Hipótesis de Trabajo}

\begin{quote}
La integración de restricciones de suavidad espacial basadas en topologías de grafos junto con \textit{priors} de continuidad temporal penalizada (tales como la regularización de Sobolev TRSS) permite reconstruir canales faltantes de EEG de manera significativamente más precisa y robusta que los métodos geométricos de extrapolación de superficies estáticas. Esta ventaja observada no solo se refleja en la mitigación del error cuadrático, sino fundamentalmente en la preservación fiel de la morfología transitoria (ERPs) y la densidad espectral (PSD) de las señales, especialmente ante escenarios clínicos de pérdida espacial severa y contigua.
\end{quote}

Al cierre de la extensa fase experimental de esta investigación, la evidencia acumulada resulta fuertemente favorable a esta hipótesis. Como se expondrá en los capítulos de resultados, la regularización temporal marca un punto de inflexión decisivo en el desempeño de interpolación, validando la eficacia del enfoque propuesto.

\section{Alcance y Aportes de la Investigación}

El alcance de este trabajo engloba la ejecución computacional a gran escala del \textit{pipeline} de evaluación cruzando múltiples datasets y escenarios de falla. A diferencia de las evaluaciones limitadas que abundan en el estado del arte, este trabajo consolida el análisis de miles de ensayos EEG bajo la optimización de decenas de miles de configuraciones paramétricas mediante Optuna.

Para garantizar reproducibilidad y cobertura sistemática, la ejecución se ancla en un registro exhaustivo de corridas experimentales que fija explícitamente las combinaciones de datasets, modos de pérdida, semillas, grafos y familias de métodos. Este diseño evita evaluaciones \textit{ad-hoc} y permite trazar cada afirmación cuantitativa reportada en los capítulos posteriores hasta su artefacto experimental correspondiente.

Los principales aportes científicos y técnicos de esta tesis pueden resumirse en tres ejes. En primer lugar, se proporciona evidencia estadística de un cruce de desempeño teórico: una vez que los hiperparámetros están debidamente sintonizados, las familias de regularización temporal (TRSS y Tikhonov) dominan de forma consistente a los \textit{baselines} geométricos en fidelidad morfológica. En segundo lugar, se revela que bajo regímenes críticos ---por ejemplo, un $40\%$ de canales faltantes agrupados localmente--- los métodos geométricos clásicos experimentan un colapso de fase, mientras que el ancla temporal de los métodos GSP mantiene la estabilidad de la trayectoria fisiológica. Finalmente, se entrega a la comunidad un \textit{toolkit} de experimentación unificado en Python, completamente rastreable y de código abierto, que permite ejecutar interpolaciones GSP avanzadas con optimización bayesiana de forma automatizada sobre datos neurológicos propios.

\section{Estructura de la Tesis}

Para desarrollar el análisis propuesto, el resto del documento se organiza de la siguiente manera. El Capítulo~2 presenta los fundamentos neurofisiológicos del EEG, los pilares matemáticos del Procesamiento de Señales en Grafos ---incluyendo el Laplaciano, la Transformada de Fourier en Grafos y la inferencia de grafos--- y el estado del arte en interpolación de señales biomédicas. El Capítulo~3 detalla el protocolo experimental: los tres datasets empleados, el modelado matemático de los algoritmos evaluados, los métodos de simulación de daños y la arquitectura del motor de optimización bayesiana. El Capítulo~4 expone los hallazgos empíricos, incluyendo comparativas globales de error, el análisis del cruce del \textit{baseline}, curvas de degradación por pérdida severa y demostraciones visuales de reconstrucción de series de tiempo y espectros. El Capítulo~5 ofrece una discusión crítica sobre las causas metodológicas de los resultados obtenidos y transparenta las limitaciones del estudio. Finalmente, el Capítulo~6 sintetiza las conclusiones, valora el cumplimiento de la hipótesis y sugiere vías de investigación futuras.

\section{Tesis, contribución y criterio de lectura}
La tesis que organiza este documento es que la reconstrucción de canales EEG faltantes debe evaluarse como un problema simultáneamente espacial, temporal y espectral. En consecuencia, una evaluación basada únicamente en MAE o RMSE no es suficiente para decidir si una interpolación es útil en aplicaciones neurofisiológicas. Esta investigación adopta una posición deliberadamente conservadora: los resultados favorables a los métodos temporales se reportan solo cuando pueden trazarse a artefactos reproducibles, y las limitaciones de equivalencia, normalización o disponibilidad de datos se explicitan en los capítulos correspondientes.

La contribución se puede sintetizar en una oración: \emph{esta tesis establece un benchmark reproducible y justo para reconstrucción de canales EEG, mostrando que la regularización temporal sobre grafos mejora la fidelidad morfológica y espectral bajo escenarios de pérdida severa, incluso después de optimizar rigurosamente los baselines clásicos}. Esta formulación evita presentar el trabajo como una propuesta algorítmica aislada; el aporte principal es el diseño experimental, la trazabilidad y la comparación controlada entre familias metodológicas.
